\section{Introduction}\label{sec:introduction}

Question answering over drug information supports patients, caregivers and pharmacists, and because an incorrect dosage or an omitted contraindication can cause direct harm, answers must be accurate and traceable to authoritative sources. Retrieval-augmented generation (RAG) meets this requirement by grounding a language model in retrieved documents~\cite{lewis2020rag}. This paper focuses on a specific setting of this problem, in which Vietnamese questions are answered from official drug labelling and cost constraints limit inference to low-cost language models.

This setting is challenging for three reasons. First, Vietnamese lacks an evaluation resource for drug information. Existing Vietnamese medical datasets address general health questions~\cite{huy2021vimq,nguyen2022spbertqa,luu2021vicoqa} rather than drug labelling. Drug questions often use brand names, colloquial product names or text without diacritics. Second, retrieval approaches are usually ranked by effectiveness alone~\cite{thakur2021beir,muennighoff2023mteb}. For a deployed system, however, embedding size and reranker throughput matter as much as ranking quality, and it is unclear which combination of sparse, dense, hybrid and reranked retrieval offers the best trade-off. Finally, agentic RAG methods either fine-tune the generator or rely on a capable model to decide when and how to retrieve~\cite{asai2024selfrag,yan2024crag,jeong2024adaptiverag}. Such pipelines are also typically reported with a single aggregate score, which reveals neither the contribution of each step nor the behaviour on unanswerable or adversarial requests.

To address these challenges, we first construct a Vietnamese drug-information corpus from the national drug formulary and package leaflets, together with a retrieval benchmark whose answer locations are known by construction. We then use this benchmark to compare retrieval approaches under a common implementation, reporting their cost next to their effectiveness. Finally, we design an agentic RAG pipeline for low-cost models. Our key design choice is to keep the control flow in deterministic code and to restrict every model call to a single, small structured decision, such as whether a request is safe or whether the retrieved evidence is sufficient. Narrow tasks suit small models better than open-ended planning, and the code enforces budgets and lets each step be ablated.

The experiments ask which retrieval approach offers the best trade-off between effectiveness and cost, how much the pipeline gains over the same language model without retrieval and over one-step RAG, which of its steps contribute to the gain, and whether it declines unanswerable, out-of-scope and adversarial requests. The pipeline is run with Gemini 3.8 Flash and with the 9-billion-parameter open-weight Qwen3.5-9B, under a protocol with established correctness, attribution and safety metrics, paired significance tests and a calibrated automatic judge. The experiments show that grounding the model in the corpus reduces hallucinated answers eighteen- to twenty-five-fold for both models, that query rephrasing enables the agent to answer the conversational follow-up questions that one-step RAG leaves almost entirely unanswered while the judge--refine loop further improves fact coverage, and that within the pipeline Qwen3.5-9B comes close to Gemini 3.8 Flash.

This paper makes three contributions. It provides a Vietnamese drug-information corpus with a 10{,}000-question retrieval benchmark whose answer locations are known. It then reports a systematic comparison of sparse, dense, hybrid and reranked retrieval in which effectiveness is set against deployment cost. Finally, it presents an agentic RAG pipeline that small open-weight models can run, together with an end-to-end evaluation protocol combining ablations, safety cases and judge calibration, that quantifies the benefits and the limits of this design. The code and the evaluation reports are released on GitHub,\footnote{\url{https://github.com/doanan293/vipharma-rag}} and the corpus, the benchmark, the golden set, the retrieval outputs and the answers of every reported run are released as a public dataset.\footnote{\url{https://www.kaggle.com/datasets/doanvanan0209/vipharma-rag}}
